% =============================================================================
% InEar Teensy Compact Protocol Specification
% =============================================================================
\documentclass[11pt,a4paper]{article}

% --- Packages ---
\usepackage[utf8]{inputenc}
\usepackage[T1]{fontenc}
\usepackage{lmodern}
\usepackage[margin=2.5cm]{geometry}
\usepackage{graphicx}
\usepackage{xcolor}
\usepackage{booktabs}
\usepackage{tabularx}
\usepackage{array}
\usepackage{multirow}
\usepackage{makecell}
\usepackage{amsmath}
\usepackage{amssymb}
\usepackage{enumitem}
\usepackage{listings}
\usepackage{float}
\usepackage{caption}
\usepackage{textcomp}
\usepackage{hyperref}
\usepackage{cleveref}
\usepackage{tikz}
\usepackage{bytefield}
\usepackage{fancyhdr}
\usepackage{titlesec}
\usepackage{pifont}

% --- Colors ---
\definecolor{backgray}{HTML}{F7F7F7}
\definecolor{lightgray}{HTML}{D0D0D0}
\definecolor{accentblue}{HTML}{2563EB}

% --- Listings setup ---
\lstdefinestyle{code}{
    backgroundcolor=\color{backgray},
    basicstyle=\ttfamily\small,
    breaklines=true,
    frame=single,
    rulecolor=\color{black},
    numbers=none,
    tabsize=4,
    showstringspaces=false,
    commentstyle=\itshape,
    keywordstyle=\bfseries,
    stringstyle={},
}

\lstdefinestyle{cpp}{
    style=code,
    language=C++,
    morekeywords={uint8_t, uint16_t, uint32_t, int32_t, int64_t, size_t},
}

\lstdefinestyle{diagram}{
    backgroundcolor=\color{backgray},
    basicstyle=\ttfamily\footnotesize,
    breaklines=false,
    frame=single,
    rulecolor=\color{black},
    numbers=none,
    tabsize=2,
    showstringspaces=false,
}

\lstset{style=code}

% --- Custom commands ---
\newcommand{\hexval}[1]{\texttt{0x#1}}
\newlength{\pktbitwidth}
\setlength{\pktbitwidth}{14.2pt}
\hfuzz=30pt % suppress bytefield overfull warnings (diagrams span full margin width)

% --- Header/Footer ---
\pagestyle{fancy}
\fancyhf{}
\fancyhead[L]{\small InEar Teensy Compact Protocol}
\fancyhead[R]{\small v1.0}
\fancyfoot[C]{\thepage}
\renewcommand{\headrulewidth}{0.4pt}

% --- Section formatting ---
\titleformat{\section}
    {\Large\bfseries}{\thesection}{1em}{}
\titleformat{\subsection}
    {\large\bfseries}{\thesubsection}{1em}{}
\titleformat{\subsubsection}
    {\normalsize\bfseries}{\thesubsubsection}{1em}{}

% --- Hyperref setup ---
\hypersetup{
    colorlinks=true,
    linkcolor=black,
    urlcolor=accentblue,
    citecolor=black,
    pdftitle={InEar Teensy Compact Protocol Specification},
    pdfauthor={InEar EEG Project},
    pdfsubject={EEG Streaming Protocol},
}

% --- Compact lists ---
\setlist{nosep, leftmargin=1.5em}

% --- Table column types ---
\newcolumntype{C}[1]{>{\centering\arraybackslash}p{#1}}
\newcolumntype{R}[1]{>{\raggedleft\arraybackslash}p{#1}}
\newcolumntype{L}[1]{>{\raggedright\arraybackslash}p{#1}}

% =============================================================================
\begin{document}

\begin{center}
    {\LARGE\bfseries InEar Teensy Compact Protocol} \\[0.5em]
    {\large Data Protocol Specification} \\[1em]
    {\small Version 1.0 \quad---\quad February 2026}
\end{center}

\vspace{1em}
\tableofcontents
\vspace{1.5em}

% =============================================================================
\newpage
\section{Protocol Overview}
\label{sec:overview}

\medskip\noindent
\textbf{Key design decisions:}

\begin{itemize}
    \item \textbf{No footer bytes} --- the 2-byte \hexval{C0}\,\hexval{C0} end marker is eliminated entirely.  The parser relies on the sync header + known packet size.
    \item \textbf{No sequence counter} --- the 4-byte microsecond timestamp provides both ordering and loss detection.
    \item \textbf{TTL merged into TYPE byte} --- two TTL bits (input + output) share a byte with the 3-bit packet type, saving the former 1-byte marker field.
    \item \textbf{EEG + at most one auxiliary sensor per packet} --- staggered scheduling avoids large combined packets.
    \item \textbf{Built-in error signaling} --- a reserved bit pattern forces the host application back into handshake mode.
\end{itemize}

\begin{table}[H]
\centering
\caption{Protocol parameters at a glance}
\label{tab:overview}
\begin{tabular}{|l|l|}
\hline
\textbf{Parameter} & \textbf{Value} \\\hline
\hline
EEG Sample Rate & 1000\,Hz \\
EEG Channels & 8 (24-bit signed, 24\,B total) \\
Baud Rate & 2,000,000 (2\,Mbaud) \\
Byte Order & Big-Endian (network order) \\
Packet Size & 32--38\,B (variable, depends on aux payload) \\
Best Case (EEG only) & 32\,B \\
Worst Case (EEG + Accel) & 38\,B \\
Fixed Overhead & 8\,B (7\,B header + 1\,B checksum) \\
Error Detection & XOR checksum (1\,B) \\
\hline
\end{tabular}
\end{table}

% =============================================================================
\section{Packet Structure}
\label{sec:packet-structure}

\begin{center}
\fbox{\texttt{START\,(2\,B)\;\textbar\; TIMESTAMP\,(4\,B)\;\textbar\; TYPE{+}TTL\,(1\,B)\;\textbar\; EEG\,(24\,B)\;\textbar\; [AUX]\;\textbar\; CHECKSUM\,(1\,B)}}
\end{center}

% ---------------------------------------------------------------------------
\subsection{Header (7 bytes --- always present)}

\begin{table}[H]
\centering
\caption{Header fields}
\label{tab:header}
\small
\begin{tabular}{|R{1.2cm}|R{0.8cm}|l|l|l|}
\hline
\textbf{Offset} & \textbf{Size} & \textbf{Field} & \textbf{Format} & \textbf{Description} \\\hline
0 & 1\,B & Sync\,1 & \hexval{A5} & Header byte 1 (constant) \\
1 & 1\,B & Sync\,2 & \hexval{5A} & Header byte 2 (constant) \\
2--5 & 4\,B & Timestamp & uint32 BE & Microseconds since boot \\
6 & 1\,B & Type{+}TTL & bitfield & Packet type + TTL lines (\cref{sec:type-ttl}) \\
\hline
\multicolumn{2}{r}{\textbf{Total:}} & \multicolumn{3}{l}{\textbf{7 bytes}} \\
\hline
\end{tabular}
\end{table}

% ---------------------------------------------------------------------------
\subsection{EEG Data (24 bytes --- always present)}

Eight channels of 24-bit signed big-endian data from the ADS1299:

\begin{table}[H]
\centering
\caption{EEG data block (8 channels $\times$ 3\,B = 24\,B)}
\label{tab:eeg-block}
\small
\begin{tabular}{|R{1.2cm}|R{0.8cm}|l|l|l|}
\hline
\textbf{Offset} & \textbf{Size} & \textbf{Field} & \textbf{Format} & \textbf{Description} \\\hline
7--9 & 3\,B & EEG Ch\,1 & int24 BE & Channel 1 \\
10--12 & 3\,B & EEG Ch\,2 & int24 BE & Channel 2 \\
13--15 & 3\,B & EEG Ch\,3 & int24 BE & Channel 3 \\
16--18 & 3\,B & EEG Ch\,4 & int24 BE & Channel 4 \\
19--21 & 3\,B & EEG Ch\,5 & int24 BE & Channel 5 \\
22--24 & 3\,B & EEG Ch\,6 & int24 BE & Channel 6 \\
25--27 & 3\,B & EEG Ch\,7 & int24 BE & Channel 7 \\
28--30 & 3\,B & EEG Ch\,8 & int24 BE & Channel 8 \\
\hline
\multicolumn{2}{r}{\textbf{Total:}} & \multicolumn{3}{l}{\textbf{24 bytes} (offsets 7--30)} \\
\hline
\end{tabular}
\end{table}

% ---------------------------------------------------------------------------
\subsection{Optional Auxiliary Block (0--6 bytes)}
\label{sec:aux-block}

At most \textbf{one} auxiliary sensor is attached per packet.  The aux block immediately follows the EEG data at offset~31.  The TYPE field (bits 2--0) determines which block, if any, is present:

\begin{table}[H]
\centering
\caption{Auxiliary data blocks by type code}
\label{tab:aux-blocks}
\small
\begin{tabular}{|c|l|R{0.8cm}|l|l|}
\hline
\textbf{Type} & \textbf{Aux Block} & \textbf{Size} & \textbf{Format} & \textbf{Sensor IC} \\\hline
\hexval{00} & (none) & 0\,B & --- & --- \\
\hexval{01} & Accelerometer & 6\,B & 3$\times$ int16 BE & ADXL367 \\
\hexval{02} & PPG + SpO\textsubscript{2} & 4\,B & uint24 BE + uint8 & MAXM86161 \\
\hexval{03} & Temperature & 2\,B & int16 BE & TMP117 \\
\hexval{04} & Battery & 1\,B & uint8 & (voltage divider) \\
\hexval{05} & Reserved\,1 & --- & --- & --- \\
\hexval{06} & Reserved\,2 & --- & --- & --- \\
\hexval{07} & Reserved\,3 & --- & --- & --- \\
\hline
\end{tabular}
\end{table}

% ---------------------------------------------------------------------------
\subsection{Checksum (1 byte --- always last)}

XOR of every preceding byte in the packet (offset~0 to $n{-}2$, where $n$ is the total packet length):

\begin{table}[H]
\centering
\small
\begin{tabular}{|R{1.2cm}|R{0.8cm}|l|l|l|}
\hline
\textbf{Offset} & \textbf{Size} & \textbf{Field} & \textbf{Format} & \textbf{Description} \\\hline
$n{-}1$ & 1\,B & Checksum & uint8 & XOR of bytes 0\,.\,.\,$(n{-}2)$ \\
\hline
\end{tabular}
\end{table}

% =============================================================================
\subsection{Packet Size}
\label{sec:packet-sizes}

\begin{table}[H]
\centering
\caption{Packet sizes for each type code}
\label{tab:packet-sizes}
\begin{tabular}{|c|l|c|c|c|}
\hline
\textbf{Type} & \textbf{Contents} & \textbf{Header} & \textbf{Aux} & \textbf{Total} \\\hline
\hline
\hexval{00} & EEG only (+ TTL) & 7 + 24 + 1 & 0\,B & \textbf{32\,B} \\
\hexval{01} & EEG + Accelerometer & 7 + 24 + 1 & 6\,B & \textbf{38\,B} \\
\hexval{02} & EEG + PPG + SpO\textsubscript{2} & 7 + 24 + 1 & 4\,B & \textbf{36\,B} \\
\hexval{03} & EEG + Temperature & 7 + 24 + 1 & 2\,B & \textbf{34\,B} \\
\hexval{04} & EEG + Battery & 7 + 24 + 1 & 1\,B & \textbf{33\,B} \\
\hexval{05} & Reserved\,1 & --- & --- & --- \\
\hexval{06} & Reserved\,2 & --- & --- & --- \\
\hexval{07} & Reserved\,3 & --- & --- & --- \\
\hline
\end{tabular}
\end{table}

% =============================================================================
\section{Packet Structure Diagrams}
\label{sec:diagrams}

\subsection{\texorpdfstring{Best Case --- EEG + TTL Only (32\,B)}{Best Case --- EEG + TTL Only (32 B)}}

\begin{figure}[H]
\centering
\begin{bytefield}[bitwidth=\pktbitwidth, bitheight=3.5em]{32}
    \bitbox{3}{\small\textbf{START}\\\scriptsize\texttt{A5\,5A}\\\scriptsize 2\,B} &
    \bitbox{5}{\small\textbf{TSMP}\\\scriptsize uint32 BE\\\scriptsize 4\,B} &
    \bitbox{3}{\small\textbf{TYPE}\\\small\textbf{+TTL}\\\scriptsize 1\,B} &
    \bitbox{18}{\normalsize\textbf{EEG} --- 8 channels $\times$ 3\,B = 24\,B} &
    \bitbox{3}{\small\textbf{CHK}\\\scriptsize XOR\\\scriptsize 1\,B}
\end{bytefield}
\caption{EEG-only packet (Type \hexval{00}) --- 32 bytes.  Most common packet (${\sim}$80\,\% of all packets).}
\label{fig:pkt-eeg-only}
\end{figure}

\subsection{\texorpdfstring{Worst Case --- EEG + Accelerometer (38\,B)}{Worst Case --- EEG + Accelerometer (38 B)}}

\begin{figure}[H]
\centering
\begin{bytefield}[bitwidth=\pktbitwidth, bitheight=3.5em]{32}
    \bitbox{3}{\small\textbf{START}\\\scriptsize\texttt{A5\,5A}\\\scriptsize 2\,B} &
    \bitbox{4}{\small\textbf{TSMP}\\\scriptsize uint32 BE\\\scriptsize 4\,B} &
    \bitbox{3}{\small\textbf{TYPE}\\\small\textbf{+TTL}\\\scriptsize 1\,B} &
    \bitbox{14}{\normalsize\textbf{EEG} --- 8 ch $\times$ 3\,B = 24\,B} &
    \bitbox{5}{\small\textbf{ACCEL}\\\scriptsize X, Y, Z\\\scriptsize 6\,B} &
    \bitbox{3}{\small\textbf{CHK}\\\scriptsize XOR\\\scriptsize 1\,B}
\end{bytefield}
\caption{EEG + Accelerometer packet (Type \hexval{01}) --- 38 bytes.  Largest packet in the protocol.}
\label{fig:pkt-eeg-accel}
\end{figure}

\subsection{\texorpdfstring{EEG + PPG + SpO\textsubscript{2} (36\,B)}{EEG + PPG + SpO2 (36 B)}}

\begin{figure}[H]
\centering
\begin{bytefield}[bitwidth=\pktbitwidth, bitheight=3.5em]{32}
    \bitbox{3}{\small\textbf{START}\\\scriptsize\texttt{A5\,5A}\\\scriptsize 2\,B} &
    \bitbox{4}{\small\textbf{TSMP}\\\scriptsize uint32 BE\\\scriptsize 4\,B} &
    \bitbox{3}{\small\textbf{TYPE}\\\small\textbf{+TTL}\\\scriptsize 1\,B} &
    \bitbox{14}{\normalsize\textbf{EEG} --- 8 ch $\times$ 3\,B = 24\,B} &
    \bitbox{3}{\small\textbf{PPG}\\\scriptsize uint24\\\scriptsize 3\,B} &
    \bitbox{2}{\small\textbf{SpO$_2$}\\\scriptsize uint8\\\scriptsize 1\,B} &
    \bitbox{3}{\small\textbf{CHK}\\\scriptsize XOR\\\scriptsize 1\,B}
\end{bytefield}
\caption{EEG + PPG + SpO\textsubscript{2} packet (Type \hexval{02}) --- 36 bytes.}
\label{fig:pkt-eeg-ppg}
\end{figure}

\subsection{\texorpdfstring{EEG + Temperature (34\,B)}{EEG + Temperature (34 B)}}

\begin{figure}[H]
\centering
\begin{bytefield}[bitwidth=\pktbitwidth, bitheight=3.5em]{32}
    \bitbox{3}{\small\textbf{START}\\\scriptsize\texttt{A5\,5A}\\\scriptsize 2\,B} &
    \bitbox{4}{\small\textbf{TSMP}\\\scriptsize uint32 BE\\\scriptsize 4\,B} &
    \bitbox{3}{\small\textbf{TYPE}\\\small\textbf{+TTL}\\\scriptsize 1\,B} &
    \bitbox{15}{\normalsize\textbf{EEG} --- 8 ch $\times$ 3\,B = 24\,B} &
    \bitbox{4}{\small\textbf{TEMP}\\\scriptsize int16 BE\\\scriptsize 2\,B} &
    \bitbox{3}{\small\textbf{CHK}\\\scriptsize XOR\\\scriptsize 1\,B}
\end{bytefield}
\caption{EEG + Temperature packet (Type \hexval{03}) --- 34 bytes.}
\label{fig:pkt-eeg-temp}
\end{figure}

\subsection{\texorpdfstring{EEG + Battery (33\,B)}{EEG + Battery (33 B)}}

\begin{figure}[H]
\centering
\begin{bytefield}[bitwidth=\pktbitwidth, bitheight=3.5em]{32}
    \bitbox{3}{\small\textbf{START}\\\scriptsize\texttt{A5\,5A}\\\scriptsize 2\,B} &
    \bitbox{4}{\small\textbf{TSMP}\\\scriptsize uint32 BE\\\scriptsize 4\,B} &
    \bitbox{3}{\small\textbf{TYPE}\\\small\textbf{+TTL}\\\scriptsize 1\,B} &
    \bitbox{15}{\normalsize\textbf{EEG} --- 8 ch $\times$ 3\,B = 24\,B} &
    \bitbox{4}{\small\textbf{BAT}\\\scriptsize uint8\\\scriptsize 1\,B} &
    \bitbox{3}{\small\textbf{CHK}\\\scriptsize XOR\\\scriptsize 1\,B}
\end{bytefield}
\caption{EEG + Battery packet (Type \hexval{04}) --- 33 bytes.}
\label{fig:pkt-eeg-bat}
\end{figure}

% =============================================================================
\newpage
\section{TYPE+TTL Byte}
\label{sec:type-ttl}

The single TYPE+TTL byte at offset~6 packs both TTL trigger lines and the packet type into 8~bits:

\begin{figure}[H]
\centering
\begin{bytefield}[bitwidth=2.5em, bitheight=3em]{8}
    \bitheader{0,1,2,3,4,5,6,7} \\
    \bitbox{1}{\small\texttt{TTL}\\\texttt{IN}} &
    \bitbox{1}{\small\texttt{TTL}\\\texttt{OUT}} &
    \bitbox{3}{\small Reserved\\\small(\texttt{000})} &
    \bitbox{3}{\small TYPE\\\small(3 bits)}
\end{bytefield}
\caption{TYPE+TTL byte bit-field layout.  Bit~7 = MSB, Bit~0 = LSB.}
\label{fig:typettl-bitfield}
\end{figure}

\subsection{Bit Assignments}

\begin{table}[H]
\centering
\caption{TYPE+TTL bit assignments}
\label{tab:typettl-bits}
\begin{tabular}{|c|l|l|}
\hline
\textbf{Bit(s)} & \textbf{Name} & \textbf{Description} \\\hline
\hline
7 & \texttt{TTL\_IN} & External trigger input (1 = active) \\
6 & \texttt{TTL\_OUT} & External trigger output (1 = active) \\
5--3 & Reserved & Must be \texttt{000} for normal packets; \texttt{111} = error signal \\
2--0 & TYPE & Packet type selector (0--7) \\
\hline
\end{tabular}
\end{table}

\subsection{Type Code Chart}

\begin{table}[H]
\centering
\caption{Type codes (bits 2--0 of TYPE+TTL byte)}
\label{tab:type-codes}
\begin{tabular}{|c|c|l|l|r|}
\hline
\textbf{Binary} & \textbf{Hex} & \textbf{Name} & \textbf{Auxiliary Data} & \textbf{Packet Size} \\\hline
\hline
\texttt{000} & \hexval{00} & \texttt{EEG} & None & 32\,B \\
\texttt{001} & \hexval{01} & \texttt{EEG\_ACCEL} & Accelerometer (6\,B) & 38\,B \\
\texttt{010} & \hexval{02} & \texttt{EEG\_PPG} & PPG + SpO\textsubscript{2} (4\,B) & 36\,B \\
\texttt{011} & \hexval{03} & \texttt{EEG\_TEMP} & Temperature (2\,B) & 34\,B \\
\texttt{100} & \hexval{04} & \texttt{EEG\_BAT} & Battery level (1\,B) & 33\,B \\
\texttt{101} & \hexval{05} & \texttt{RESERVED\_1} & Future use & --- \\
\texttt{110} & \hexval{06} & \texttt{RESERVED\_2} & Future use & --- \\
\texttt{111} & \hexval{07} & \texttt{RESERVED\_3} & Future use & --- \\
\hline
\end{tabular}
\end{table}

\subsection{TTL Examples}

\begin{table}[H]
\centering
\caption{Example TYPE+TTL byte values}
\label{tab:typettl-examples}
\small
\begin{tabular}{|c|c|c|c|l|}
\hline
\textbf{Hex} & \textbf{Binary} & \textbf{TTL} & \textbf{Type} & \textbf{Meaning} \\\hline
\hline
\hexval{00} & \texttt{0\,0\,|\,000\,|\,000} & --- & EEG & EEG only, no trigger \\
\hexval{80} & \texttt{1\,0\,|\,000\,|\,000} & IN & EEG & EEG only, TTL input active \\
\hexval{C1} & \texttt{1\,1\,|\,000\,|\,001} & IN+OUT & ACCEL & EEG + Accel, both triggers \\
\hexval{02} & \texttt{0\,0\,|\,000\,|\,010} & --- & PPG & EEG + PPG + SpO\textsubscript{2}, no trigger \\
\hexval{43} & \texttt{0\,1\,|\,000\,|\,011} & OUT & TEMP & EEG + Temp, TTL output active \\
\hexval{84} & \texttt{1\,0\,|\,000\,|\,100} & IN & BAT & EEG + Battery, TTL input active \\
\hline
\end{tabular}
\end{table}

\newpage
\subsection{Error Signaling}

When the reserved bits (5--3) are all set to \texttt{1}, the packet signals an \textbf{error condition}:

\[
\text{Error byte} = \texttt{??}\underbrace{\texttt{111}}_{\text{error flag}}\texttt{???}_{\text{binary}} \quad\longrightarrow\quad \text{bits 5--3} = \texttt{111}
\]

An error packet consists of the normal header structure with the error flag set and zeroed-out EEG data:

\begin{center}
\texttt{A5\,5A + TIMESTAMP\,(4\,B) + Error-flagged TYPE+TTL + 24$\times$\hexval{00} + CHECKSUM}
\end{center}

Upon receiving a packet with bits 5--3 = \texttt{111}, the Qt host application must:
\begin{enumerate}
    \item Stop data acquisition immediately.
    \item Transition back to the handshake / connection state.
    \item Report the error condition to the user.
\end{enumerate}

% =============================================================================
\section{Byte Offset Reference}
\label{sec:byte-offsets}

\begin{table}[H]
\centering
\caption{Absolute byte offsets for all active packet types (8-channel EEG)}
\label{tab:offsets}
\small
\begin{tabular}{|l|C{1.5cm}|C{1.5cm}|C{1.5cm}|C{1.5cm}|C{1.5cm}|}
\hline
\textbf{Block} & \textbf{\hexval{00}} & \textbf{\hexval{01}} & \textbf{\hexval{02}} & \textbf{\hexval{03}} & \textbf{\hexval{04}} \\
& \tiny EEG & \tiny +Accel & \tiny +PPG & \tiny +Temp & \tiny +Bat \\
\hline
Sync\,(2\,B)       & 0--1   & 0--1   & 0--1   & 0--1   & 0--1 \\
Timestamp\,(4\,B)  & 2--5   & 2--5   & 2--5   & 2--5   & 2--5 \\
Type+TTL\,(1\,B)   & 6      & 6      & 6      & 6      & 6 \\
EEG\,(24\,B)       & 7--30  & 7--30  & 7--30  & 7--30  & 7--30 \\
Aux data           & ---    & 31--36 & 31--34 & 31--32 & 31 \\
Checksum\,(1\,B)   & 31     & 37     & 35     & 33     & 32 \\
\hline
\textbf{Total}     & \textbf{32\,B} & \textbf{38\,B} & \textbf{36\,B} & \textbf{34\,B} & \textbf{33\,B} \\
\hline
\end{tabular}
\end{table}

% =============================================================================
\section{Multi-Rate Sensor Scheduling}
\label{sec:scheduling}

Every packet contains \textbf{fresh} EEG data sampled at 1000\,Hz.  Each auxiliary sensor has its own dedicated packet type and sub-rate.  Because the protocol allows at most one auxiliary block per packet, collisions between sensors due at the same sample are resolved by \textbf{deferring} the lower-priority sensor by 1\,ms.  The deferred packet still carries a fresh EEG sample from that later time slot --- no EEG data is ever stale or duplicated.

\begin{table}[H]
\centering
\caption{Sensor scheduling rates}
\label{tab:schedule}
\begin{tabular}{|l|r|r|}
\hline
\textbf{Sensor} & \textbf{Rate (Hz)} & \textbf{Every $N$th pkt}\\\hline
\hline
EEG (+ TTL) & 1000 & Every packet\\
Accelerometer & 200 & Every 5th\\
PPG + SpO\textsubscript{2} & 100 & Every 10th\\
Temperature & 1 & Every 1000th\\
Battery & 1 & Every 1000th\\
\hline
\end{tabular}
\end{table}

\subsection{\texorpdfstring{Collision Resolution: 1\,ms Deferral}{Collision Resolution: 1 ms Deferral}}

When two or more auxiliary sensors are due on the same sample, the highest-priority sensor claims that packet and lower-priority sensors are deferred by 1\,ms (one sample period).
\medskip\noindent
\textbf{Priority order (highest first):}
\[
\text{ACCEL} \;>\; \text{PPG} \;>\; \text{TEMP} \;>\; \text{BAT}
\]

\subsection{\texorpdfstring{1\,Hz Sensor Staggering}{1 Hz Sensor Staggering}}

Temperature and Battery are both sampled at 1\,Hz (every 1000th packet).  To avoid them always colliding, they are staggered with a fixed offset so that only one 1\,Hz sensor is due at any given second boundary.

\medskip\noindent
\textbf{Note:} The exact staggering offset between Temperature and Battery is \textbf{TBD} (not yet finalized in the firmware).  The only hard constraint is:

\begin{itemize}
    \item Temperature and Battery must not be due on the same sample.
    \item If either 1\,Hz sensor collides with an Accel or PPG packet, the 1\,Hz sensor is deferred by 1\,ms (same deferral rule as above).
\end{itemize}

\subsection{Example: 21 Consecutive Packets}

\begin{lstlisting}[style=diagram, caption={Scheduling over 21 samples (Accel every 5th, PPG every 10th)}]
Pkt#:   1     2     3     4     5     6     7     8     9     10
EEG:   [EEG] [EEG] [EEG] [EEG] [EEG] [EEG] [EEG] [EEG] [EEG] [EEG]
Aux:    ---   ---   ---   ---  ACCEL  ---   ---   ---   ---  ACCEL
Size:   32B   32B   32B   32B   38B   32B   32B   32B   32B   38B

Pkt#:   11    12    13    14    15    16    17    18    19    20    21
EEG:   [EEG] [EEG] [EEG] [EEG] [EEG] [EEG] [EEG] [EEG] [EEG] [EEG] [EEG]
Aux:    PPG   ---   ---   ---  ACCEL  ---   ---   ---   ---  ACCEL  PPG
Size:   36B   32B   32B   32B   38B   32B   32B   32B   32B   38B   36B

Note: PPG due at pkt 10 deferred to pkt 11 (Accel had priority).
      PPG due at pkt 20 deferred to pkt 21 (Accel had priority).
      Each deferred packet carries fresh EEG from its own time slot.
\end{lstlisting}


% =============================================================================
\newpage
\section{LSL Stream Metadata Reference}
\label{sec:lsl}

\subsection{StreamInfo Properties}

\begin{table}[H]
\centering
\caption{LSL \texttt{StreamInfo} constructor parameters}
\label{tab:lsl-streaminfo}
\small
\begin{tabularx}{\textwidth}{|l|l|l|>{\raggedright\arraybackslash}X|}
\hline
\textbf{Property} & \textbf{EEG Stream} & \textbf{Marker Stream} & \textbf{Notes} \\\hline
\hline
\texttt{name} &
\texttt{"InEarEEG"} &
\texttt{"InEarEEG\_Markers"} &
Human-readable name.  Must be unique on the network if multiple devices are present. \\
\hline
\texttt{type} &
\texttt{"EEG"} &
\texttt{"Markers"} &
Content type.  \texttt{"EEG"} enables automatic recognition by LabRecorder and BIDS tools. \\
\hline
\texttt{channel\_count} &
\texttt{8} &
\texttt{1} &
Must match the actual channel count.  Marker streams always use 1 channel. \\
\hline
\texttt{nominal\_srate} &
\texttt{1000.0} &
\texttt{0} &
Sample rate in Hz.  Use \texttt{0} (\texttt{IRREGULAR\_RATE}) for event streams. \\
\hline
\texttt{channel\_format} &
\texttt{cf\_float32} &
\texttt{cf\_string} &
\texttt{float32} for EEG data in $\mu$V; \texttt{string} for human-readable event labels. \\
\hline
\texttt{source\_id} &
\texttt{"inear\_001"} &
\texttt{"inear\_001\_markers"} &
Unique per-device identifier.  Should remain stable across sessions. \\
\hline
\end{tabularx}
\end{table}

\subsection{XML Channel Metadata}

The \texttt{desc()} XML tree embeds per-channel metadata that downstream tools use for labeling, unit conversion, and montage reconstruction:

\begin{lstlisting}[style=code, language=XML, caption={Recommended LSL XML metadata for EEG stream}]
<channels>
  <channel>
    <label>EEG1</label>
    <unit>microvolts</unit>
    <type>EEG</type>
  </channel>
  <channel>
    <label>EEG2</label>
    <unit>microvolts</unit>
    <type>EEG</type>
  </channel>
  <!-- ... channels 3-8 ... -->
</channels>
<acquisition>
  <manufacturer>InEar Teensy</manufacturer>
  <model>Compact Protocol v1</model>
  <serial_number>DEVICE_SERIAL</serial_number>
  <scale_factor>0.02235</scale_factor>
</acquisition>
\end{lstlisting}

\subsection{Important Metadata Fields}

\begin{table}[H]
\centering
\caption{LSL XML metadata fields that matter for interoperability}
\label{tab:lsl-xml-fields}
\small
\begin{tabularx}{\textwidth}{|l|l|l|>{\raggedright\arraybackslash}X|}
\hline
\textbf{XML Path} & \textbf{Key} & \textbf{Example} & \textbf{Why It Matters} \\\hline
\hline
\texttt{channels/channel} & \texttt{label} & \texttt{Fp1}, \texttt{Cz} &
Channel names for montage.  10--20 system names enable automatic montage in MNE-Python and EEGLAB. \\
\hline
\texttt{channels/channel} & \texttt{unit} & \texttt{microvolts} &
Physical unit.  Must be \texttt{"microvolts"} (not \texttt{"uV"}) for consistent cross-tool behavior.\footnote{Observed codebase inconsistency: Python scripts use \texttt{"microvolts"} while C++ plugins use \texttt{"uV"}.  Standardize on \texttt{"microvolts"} to match the LSL/XDF conventions used by LabRecorder and MNE-Python.} \\
\hline
\texttt{channels/channel} & \texttt{type} & \texttt{EEG} &
Per-channel type.  Enables mixed-type streams (e.g., channels 1--8 = EEG). \\
\hline
\texttt{acquisition} & \texttt{manufacturer} & \texttt{InEar Teensy} &
Identifies hardware.  Used for provenance tracking in BIDS datasets. \\
\hline
\texttt{acquisition} & \texttt{model} & \texttt{Compact v1} &
Protocol / firmware version.  Helps distinguish recordings made with different firmware. \\
\hline
\texttt{acquisition} & \texttt{serial\_number} & \texttt{INEAR\_042} &
Hardware serial.  Enables per-device calibration and multi-device setups. \\
\hline
\texttt{acquisition} & \texttt{scale\_factor} & \texttt{0.02235} &
$\mu$V per LSB.  If the stream pushes raw ADC counts, tools need this to convert to microvolts. \\
\hline
\end{tabularx}
\end{table}

\subsection{Marker Stream Configuration}

\begin{table}[H]
\centering
\caption{Marker stream properties}
\label{tab:lsl-markers}
\begin{tabular}{|l|l|l|}
\hline
\textbf{Property} & \textbf{Value} & \textbf{Rationale} \\\hline
\hline
\texttt{type} & \texttt{"Markers"} & Standard LSL type for event streams \\
\texttt{nominal\_srate} & \texttt{0} & Irregular rate --- events arrive asynchronously \\
\texttt{channel\_format} & \texttt{cf\_string} & Human-readable labels \\
\texttt{channel\_count} & \texttt{1} & Single event channel \\
\hline
\end{tabular}
\end{table}

\begin{table}[H]
\centering
\caption{Marker string formats}
\label{tab:lsl-marker-strings}
\begin{tabular}{|l|l|}
\hline
\textbf{Event} & \textbf{Marker String} \\\hline
\hline
TTL input rising edge & \texttt{"TTL\_IN\_1"} \\
TTL input falling edge & \texttt{"TTL\_IN\_0"} \\
TTL output rising edge & \texttt{"TTL\_OUT\_1"} \\
TTL output falling edge & \texttt{"TTL\_OUT\_0"} \\
Experiment start & \texttt{"EXPERIMENT\_START"} \\
Experiment stop & \texttt{"EXPERIMENT\_STOP"} \\
\hline
\end{tabular}
\end{table}

\subsection{Outlet Configuration}

\begin{table}[H]
\centering
\caption{LSL \texttt{StreamOutlet} parameters}
\label{tab:lsl-outlet}
\begin{tabular}{|l|l|l|}
\hline
\textbf{Parameter} & \textbf{Value} & \textbf{Rationale} \\\hline
\hline
\texttt{chunk\_size} & $\lfloor \text{srate} / 20 \rfloor$ (= 50) & ${\sim}$50\,ms chunks for latency / efficiency balance \\
\texttt{max\_buffered} & 360 (seconds) & 6-minute buffer for brief network stalls \\
\hline
\end{tabular}
\end{table}

\end{document}
